\documentclass[10pt,a4paper]{amsart}
\usepackage[margin=19mm]{geometry}
\usepackage{amsmath,amssymb,mathtools}
\usepackage[scr=rsfs,cal=euler]{mathalfa}
\usepackage{xfrac}
\usepackage[unicode,hidelinks,bookmarks=false]{hyperref}
\newcommand{\ie}{i.e.\ }
\newcommand{\cf}{cf.\ }
\newcommand{\resp}{respectively\ }
\newcommand{\Subsection}{\subsection}
\providecommand{\nobreakdash}{\nobreak\hbox{-}}
\setlength{\emergencystretch}{2em}
\multlinegap0pt
\usepackage{amssymb}
\usepackage{enumerate}
\usepackage{xcolor}
\newtheorem{thm}{Theorem}
\DeclareMathOperator{\HS}{HS}
\DeclareMathOperator{\tr}{tr}
\title{Entropic optimal transport beyond product reference couplings:\\ the Gaussian case on Euclidean space}
\author{Paul Freulon, Nikitas Georgakis, Victor Panaretos}
\date{Source: arXiv:2507.01709v2 (CC BY 4.0)}
\begin{document}
\maketitle

\noindent\emph{Source and licence.} Entropic optimal transport beyond product reference couplings:\\ the Gaussian case on Euclidean space. Version arXiv:2507.01709v2 (CC BY 4.0), distributed by arXiv under the Creative Commons Attribution 4.0 International licence, http://creativecommons.org/licenses/by/4.0/, as recorded in the arXiv licence metadata for this exact version and shown on the abstract page https://arxiv.org/abs/2507.01709v2. Full licence notice: \texttt{licenses/arxiv-2507.01709-CC-BY-4.0.txt}.\par\smallskip

\noindent\textit{Editorial scope.} Counted unit: the article's thm thm:optim\_potentials (source0.tex lines 845--857). The article's complete original proof of it is delivered with it (source0.tex lines 1602--1660, one \texttt{proof} environment of 59 lines, which the article places away from the statement). No mathematics is written, repaired or re-derived: the statement and the proof are the article's own lines, in the article's own notation, and no retained line is substituted or re-flowed.  Retained uncounted as the source-local context the proof needs: lines 782--788.  Nothing was omitted from any retained span, so the package carries no omission record.  No retained line contains a citation, so the wrapper carries no bibliography and no bibliography entry.  No retained line contains a figure environment, an image-inclusion command or drawing code, so no image is delivered and the package declares no asset dependency.  Editorial formatting only, in addition: an AMS \texttt{amsart} preamble that restates the article's own declarations and macros used by the retained lines, namely the environments \texttt{thm} on separate counters, together with the standard packages those lines need.  The retained environments print in this document as Theorem 1 in order of appearance, whereas the article prints the same items with the numbering of its own full document.  The complete native source of the article as distributed in the arXiv e-print for this exact version is bundled unchanged as \texttt{sources/180899/source0.tex}.
    \begin{align}\label{eq:duality_OTGauss}
        W_\Sigma^\varepsilon(\mu, \nu) =  \max_{(F,G)}  ~& \left\{\langle I_d-F, A \rangle_{\HS}  + \langle I_d -G, B \rangle_{\HS} 
          + \varepsilon \log  \det \begin{pmatrix}
         F    & -M_{\varepsilon}    \\
    -M_{\varepsilon}^T&    G 
    \end{pmatrix} \right\}\\ & \quad +  \varepsilon \left[ \langle \Gamma_{11}, A \rangle_{\HS} + \langle \Gamma_{22}, B \rangle_{\HS} -\log \det(\varepsilon \Sigma^{-1})\right], \nonumber
    \end{align}

\begin{thm}\label{thm:optim_potentials}
 If $\varepsilon >0$, the optimal dual variables $F_\varepsilon^{\star}, G_\varepsilon^{\star}$ associated to dual problem \eqref{eq:duality_OTGauss} are given by the formulae
 \begin{equation}\label{eq:optim_dual}
     \begin{array}{ccc}
        F_\varepsilon^{\star} &= &  A^{-1}\left( \frac{\varepsilon}{2}I_d + \left(AM_{\varepsilon}BM_{\varepsilon}^T+ \frac{\varepsilon^2}{4}I_d  \right)^{1/2} \right)\\
        G_\varepsilon^{\star} &=  & B^{-1}\left( \frac{\varepsilon}{2}I_d + \left(BM_\varepsilon^T AM_{\varepsilon}+ \frac{\varepsilon^2}{4}I_d  \right)^{1/2} \right).
     \end{array}
 \end{equation}
We can also express the second optimal dual variable $G_\varepsilon^\star$ as a function of $F_\varepsilon^\star$ through the relation
 \begin{equation}
      G_\varepsilon^{\star} =  \varepsilon B^{-1}+ M_\varepsilon^{T}(F_\varepsilon^{\star})^{-1}M_\varepsilon.
 \end{equation}
\end{thm}

\begin{proof}
	Using the notations from Theorem \ref{thm:optim_potentials}, we can write the optimal transport cost between $\mu$ and $\nu$ as
	\begin{equation*}
		W_\Sigma^{\varepsilon}(\mu, \nu) = \langle I_d+\varepsilon \Gamma_{11} - F_{\varepsilon}^{\star}, A \rangle_{\HS} + \langle I_d + \varepsilon \Gamma_{22} - G_\varepsilon^{\star}, B \rangle_{\HS} + \varepsilon \log \det \begin{pmatrix}
			F_\varepsilon^{\star}  & -M_{\varepsilon}    \\
			-M_{\varepsilon}^T&    G_\varepsilon^{\star}
		\end{pmatrix} - \varepsilon \log \det(\varepsilon \Sigma^{-1}).
	\end{equation*}
	For the scalar product terms, we begin by writing 
	\begin{align*}
		\langle I_d +\varepsilon \Gamma_{11} - F_{\varepsilon}^{\star}, A \rangle_{\HS} & = \tr(A) + \varepsilon \langle \Gamma_{11}, A \rangle_{\HS} - \tr\left(\frac{\varepsilon}{2}I_d+ \left(AM_{\varepsilon}BM_{\varepsilon}^T+ \frac{\varepsilon^2}{4}I_d \right)^{1/2} \right) \\
		&  =  \tr(A) + \varepsilon \langle \Gamma_{11}, A \rangle_{\HS} - \tr\left( \left(AM_{\varepsilon}BM_{\varepsilon}^T+ \frac{\varepsilon^2}{4}I_d \right)^{1/2} \right) - \varepsilon \frac{d}{2}.
	\end{align*}
	And second we write, 
	\begin{align*}
		\langle I_d+\varepsilon \Gamma_{22}- G_{\varepsilon}^{\star}, B \rangle_{\HS} & = \tr(B) + \varepsilon \langle \Gamma_{22}, B \rangle_{\HS} - \tr\left(\frac{\varepsilon}{2}I_d +\left(BM_{\varepsilon}^T AM_{\varepsilon}+ \frac{\varepsilon^2}{4}I_d \right)^{1/2} \right) \\
		& = \tr(B) + \varepsilon \langle \Gamma_{22}, B \rangle_{\HS} - \tr\left(\left(BM_\varepsilon^T AM_{\varepsilon}+ \frac{\varepsilon^2}{4}I_d \right)^{1/2} \right) - \varepsilon \frac{d}{2}.
	\end{align*}
	Then, the identity
	\begin{equation*}
		M^{-1}A^{-1}\left(AM_{\varepsilon}BM_{\varepsilon}^T+ \frac{\varepsilon^2}{4}I_d  \right)^{1/2} AM_{\varepsilon}=\left(BM_{\varepsilon}^T AM_{\varepsilon}+ \frac{\varepsilon^2}{4}I_d  \right)^{1/2}
	\end{equation*}
	ensures that 
	\begin{equation*}
		\tr\left(\left(BM_{\varepsilon}^T AM_{\varepsilon}+ \frac{\varepsilon^2}{4}I_d \right)^{1/2} \right) = \tr\left(\left(AM_{\varepsilon}BM_{\varepsilon}^T+ \frac{\varepsilon^2}{4}I_d \right)^{1/2} \right).
	\end{equation*}
	We thus have that 
	\begin{align*}
		\langle I_d+\varepsilon \Gamma_{11} - F_{\varepsilon}^\star, A \rangle_{\HS} + \langle I_d + \varepsilon \Gamma_{22} - G_\varepsilon^{\star}, B \rangle_{\HS} & = \tr(A) + \tr(B) +\varepsilon \langle \Gamma_{11}, A \rangle_{\HS}  +\varepsilon \langle \Gamma_{22}, B \rangle_{\HS} \\
		&\quad  - 2 \tr\left( \left(AM_{\varepsilon}BM_{\varepsilon}^T+ \frac{\varepsilon^2}{4}I_d \right)^{1/2} \right) - \varepsilon d.
	\end{align*}
	
	
	Regarding the determinant term, we write it as 
	\begin{align*}
		\det \begin{pmatrix}
			F_\varepsilon^\star   & -M_{\varepsilon}    \\
			-M_{\varepsilon}^T&    G_\varepsilon^\star
		\end{pmatrix}& =  \det \begin{pmatrix}
			F_\varepsilon^\star  & -M_{\varepsilon}    \\
			-M_{\varepsilon}^T&    \varepsilon B^{-1}+ M_\varepsilon^{T}(F_\varepsilon^{\star})^{-1}M_\varepsilon
		\end{pmatrix} \\
		& = \det ( F_\varepsilon^{\star}) \det(\varepsilon B^{-1}+ M_\varepsilon^{T}(F_\varepsilon^{\star})^{-1}M_{\varepsilon}- M_\varepsilon^{T}(F_\varepsilon^{\star})^{-1}M) \\
		& = \varepsilon^d\det(A)^{-1}\det(B)^{-1} \det\left( \frac{\varepsilon}{2}I_d + \left(AM_{\varepsilon}BM_{\varepsilon}^T+ \frac{\varepsilon^2}{4}I_d \right)^{1/2} \right).
	\end{align*}
	Taking the logarithm of the determinant yields
	\begin{equation*}
		\varepsilon\log \det\begin{pmatrix}
			F_\varepsilon^\star   & -M_{\varepsilon}    \\
			-M_{\varepsilon}^T&    G_\varepsilon^\star
		\end{pmatrix} =  \varepsilon d \log(\varepsilon) - \varepsilon\log \det(AB) + \varepsilon \log \det\left( \frac{\varepsilon}{2}I_d + \left(AM_{\varepsilon}BM_{\varepsilon}^T+ \frac{\varepsilon^2}{4}I_d \right)^{1/2} \right).
	\end{equation*}
	
	Recalling the additive constant $-\varepsilon (2d \log(\varepsilon)+\log \det(\Sigma^{-1})$, and collecting the pieces we derive
	\begin{align*}
		W_{\Sigma}^{\varepsilon}(\mu, \nu)& = \tr(A) + \tr(B) - 2 \tr\left( \left(AM_{\varepsilon}BM_{\varepsilon}^T+ \frac{\varepsilon^2}{4}I_d \right)^{1/2} \right) + \varepsilon \log \det\left( \frac{\varepsilon}{2}I_d + \left(AM_{\varepsilon}BM_{\varepsilon}^T+ \frac{\varepsilon^2}{4}I_d \right)^{1/2} \right) \\
		&+\varepsilon \tr (\Gamma_{11} A )  +\varepsilon \tr (\Gamma_{22} B ) - \varepsilon \log \det(AB) - \varepsilon d - \varepsilon d \log(\varepsilon) - \varepsilon\log \det(\Sigma^{-1}).
	\end{align*}
\end{proof}

\end{document}
